\chapter*{서문: 우리가 알고 있는 것과 아직 모르는 것}
\addcontentsline{toc}{chapter}{서문: 우리가 알고 있는 것과 아직 모르는 것}

\begin{factbox}
가중치, 입력, 실행 코드, 수치 정밀도와 난수 상태가 고정되면 decoder-only Transformer의 순전파에서 계산되는 모든 텐서는 원칙적으로 재현할 수 있다.
그러나 그 텐서 계산이 학습한 인간 이해 가능한 특징, 회로, 알고리즘을 완전하게 기술하는 일반 이론은 아직 없다.
\end{factbox}

이 책의 목표는 이 두 문장을 동시에 이해하게 만드는 것이다.
첫 번째 문장은 Transformer를 신비화하지 않는다.
모델은 명시된 선형변환, 정규화, 비선형함수, attention, 잔차 덧셈을 실행한다.
두 번째 문장은 현재 과학의 경계를 과장하지 않는다.
행렬곱을 모두 기록하는 것과 그 계산을 설명하는 것은 동일하지 않다.

이 책의 중심 질문은 순전파의 정확한 계산과 학습된 메커니즘의 인과적 설명을 연결하는 것이다.

\begin{quote}
문자열이 입력된 순간부터 다음 토큰이 선택될 때까지 텐서의 모양과 값은 어떻게 변하는가?
그리고 학습된 계산에서 어떤 표현, 정보 흐름, 회로와 알고리즘을 인과적으로 확인할 수 있는가?
\end{quote}

첫 질문에는 정확한 수식과 실행 가능한 코드로 답할 수 있다.
두 번째 질문에는 모델·과제·입력 분포·개입 방법을 제한한 실험적 답만 줄 수 있는 경우가 많다.
이 책은 두 종류의 답을 문장 수준에서 구분한다.

\section*{범위}

주된 분석 대상은 자기회귀적 decoder-only Transformer이다.
특히 공개 논문과 구현을 함께 확인할 수 있는 GPT 계열, Llama 계열과 소형 연구 모델을 중심으로 한다.
원래의 encoder--decoder Transformer는 역사와 수식의 출발점으로 다루되, 현대 언어 생성에 필요한 causal masking, pre-normalization, RMSNorm, RoPE, SwiGLU, GQA, KV cache, prefill/decode를 별도로 전개한다.

훈련은 순전파가 왜 그런 가중치를 갖게 되는지 이해하는 데 필요한 범위에서 다룬다.
next-token likelihood, 역전파, AdamW, scaling law와 대표적인 post-training 목표를 설명한다.
분산 학습 시스템, 데이터 거버넌스, 멀티모달 모델, 검색·도구 사용 시스템은 핵심 범위 밖이다.

\section*{이 책이 약속하지 않는 것}

이 원고는 ``모든 LLM이 같은 내부 알고리즘을 사용한다''고 주장하지 않는다.
특정 회로 연구의 결론을 그 연구가 조사하지 않은 모델과 입력으로 확대하지 않는다.
또한 attention heatmap, 높은 probe 정확도, 성공한 activation steering 가운데 어느 것도 그 자체로 완전한 기계론적 설명이라고 부르지 않는다.

출판 가능한 원고라는 목표는 오류가 없다는 선언이 아니라 다음의 검증 가능성을 뜻한다.

\begin{itemize}
  \item 독자가 수식의 차원과 연산을 독립적으로 확인할 수 있다.
  \item 경험적 주장에는 원 연구와 적용 범위가 붙는다.
  \item 관찰, 인과적 증거, 해석과 열린 문제를 구분한다.
  \item 공개 모델에서 핵심 실험을 재현할 수 있는 절차가 있다.
  \item 기준일 이후의 연구로 수정될 수 있는 부분을 명시한다.
\end{itemize}

\section*{읽는 방법}

Phase 1--3은 계산을 이해하는 구간이다.
이 구간에서는 수식을 손으로 유도하고 작은 텐서로 직접 계산해야 한다.
Phase 4는 내부 상태에서 정보를 읽는 방법과 그 한계를 배운다.
Phase 5--6은 구성요소를 회로로 묶고 인과적 개입으로 가설을 시험하는 구간이다.
Phase 7은 이 도구를 factual recall, in-context learning, reasoning, refusal, hallucination 같은 행동에 적용한 연구를 비판적으로 읽는다.

각 Phase의 마지막 학습 점검은 단순 복습 문제가 아니다.
정의를 말할 수 있는지보다 수식, 텐서, 코드, 개입 결과를 서로 번역할 수 있는지를 검사한다.
이 번역 능력이 내부 메커니즘 연구의 실제 기초다.

\section*{판본과 갱신}

이 원고의 문헌 기준일은 2026년 8월 23일이다.
빠르게 변하는 전면부 연구는 날짜와 모델 버전을 함께 기록한다.
peer review를 거치지 않은 preprint와 연구기관의 기술 보고서는 그 상태를 숨기지 않는다.
새로운 반증이나 더 강한 재현 결과가 나오면 본문의 서술 강도도 함께 낮추거나 높여야 한다.
